
Machine learning research involves iterative hypothesis testing constrained by computational resources. Researchers implement hypotheses, measure performance and overhead, and discover viability constraints post-implementation. This pattern---investing hours to days before identifying critical feasibility issues---represents a methodological gap that this work addresses.

The framework validation documented in this research directory originated from a constraint-driven research question: how to design testable ML hypotheses using only existing datasets and benchmarks, without new data collection or human evaluation. From this constraint emerged a specific failure case (designated h-e1) where layer-wise logit extraction for model analysis incurred 68.65\% computational overhead, rendering the approach infeasible for real-time deployment despite technical correctness. A 10-sample micro-pilot would have revealed the overhead early, enabling an informed decision before full implementation.

This example reflects a broader methodological need. ML research workflows include hypothesis generation and performance optimization but lack formalized viability gates---systematic checkpoints that answer ''should we implement this approach at all?'' before resource commitment. Ablation studies test variations (which dropout rate performs best?) but assume viability. Big-O complexity analysis provides theoretical bounds but misses constant factors and hardware specifics that dominate real-world constraints. Expert intuition is informal and unvalidated. Full implementation provides ground truth but wastes effort on infeasible hypotheses.

The framework presented here treats viability assessment as incremental empirical validation. The core mechanism is linear overhead scaling: computational overhead measured at 10-sample micro-pilot scale predicts full-dataset overhead via extrapolation. This observation enables Gate 1 decisions before significant resource investment. Bayesian refinement at Gate 2 (100 samples) reduces prediction uncertainty when initial estimates are borderline.

Validation proceeded through a dependency chain of four hypotheses: H-E1 (corpus existence), H-M1 (micro-pilot correlation), H-M2 (Bayesian error reduction), and H-M3 (Gate 1 classification accuracy). All four passed their success criteria with statistical significance. However, the corpus was synthetic rather than drawn from published papers, and perfect linear scaling (r=1.000) represents an artifact unlikely to replicate in real data.

This work contributes:

\begin{enumerate}
\item A formalized framework for early viability assessment combining fail-fast gating with Bayesian uncertainty reduction.
\end{enumerate}

\begin{enumerate}
\item Empirical validation on synthetic data demonstrating 93.3\% Gate 1 accuracy, perfect linear correlation (r=1.000), and 40.91\% Bayesian error reduction.
\end{enumerate}

\begin{enumerate}
\item Transparent documentation of limitations: synthetic corpus, single threshold tested (10\%), untested user compliance, and unknown multi-threshold behavior.
\end{enumerate}

\begin{enumerate}
\item Prioritized future work grounded in observed limitations: real corpus validation (test $r\geq 0.7$ on published papers), prospective user study (measure stop rate and time savings), and multi-threshold validation (establish applicability scope).
\end{enumerate}

The remainder of this paper is structured as follows. Section 2 reviews related work in ablation studies, complexity analysis, early stopping, and Bayesian optimization. Section 3 describes the framework methodology including linear scaling extrapolation, Bayesian posterior refinement, and threshold-based classification. Section 4 details the experimental setup using a synthetic 32-hypothesis corpus. Section 5 presents validation results for the three mechanism hypotheses (H-M1, H-M2, H-M3). Section 6 discusses findings, limitations, and broader implications. Section 7 concludes with future directions.
